\documentclass[11pt]{article}

\usepackage[a4paper,margin=1in]{geometry}
\usepackage{amsmath,amssymb,amsthm,mathtools}
\usepackage{graphicx}
\usepackage{float}
\usepackage{microtype}
\usepackage[hidelinks]{hyperref}

\newtheorem{theorem}{Theorem}
\newtheorem{proposition}[theorem]{Proposition}
\newtheorem{corollary}[theorem]{Corollary}
\newtheorem{remark}[theorem]{Remark}

\newcommand{\C}{\mathbb C}
\newcommand{\R}{\mathbb R}
\newcommand{\N}{\mathbb N}
\newcommand{\Var}{\operatorname{Var}}
\newcommand{\dmu}{\,d\mu_{a,\beta}}
\newcommand{\cL}{\mathcal L}

\title{The exact Poincar\'e constant of the two-particle\\
beta-Ginibre ensemble}
\author{Candidate exact solved-subcase packet}
\date{9 August 2026}

\begin{document}
\maketitle

\begin{center}
\textbf{Status: candidate exact solved subcase, likely valid pending expert review.}
\end{center}

\begin{abstract}
Chafa\"i and Lehec ask for the Poincar\'e and logarithmic Sobolev constants
of the beta-Ginibre ensemble.  We determine the optimal Poincar\'e constant
for two particles and every repulsion exponent $\beta\geq0$.  For the law
on $\C^2$ with density proportional to
\[
 e^{-a(|z_1|^2+|z_2|^2)}|z_1-z_2|^\beta,\qquad a>0,
\]
the spectral gap is
\[
 \lambda_{a,\beta}=a\bigl(\sqrt{\beta^2+4}-\beta\bigr),
\]
and the optimal Poincar\'e constant is its reciprocal,
$(\sqrt{\beta^2+4}+\beta)/(4a)$.  The proof factors center of mass from
relative position, then diagonalizes the relative generator by angular
Fourier modes and generalized Laguerre polynomials.  For the normalization
of arXiv:1805.00708 and two particles, $a=2$.
\end{abstract}

\section{Source question and the result}

The source is D.~Chafa\"i and J.~Lehec, \emph{On Poincar\'e and logarithmic
Sobolev inequalities for a class of singular Gibbs measures},
arXiv:1805.00708 \cite{ChafaiLehec}.  In its comments and open questions it
considers the beta-Ginibre density
\[
 z\in\C^n\longmapsto
 e^{-n\sum_{j=1}^n|z_j|^2}\prod_{j<k}|z_j-z_k|^\beta
\]
The exact source sentence is: ``it is not clear at all what are the
Poincar\'e and log-Sobolev constants of the Ginibre ensemble.''

\begin{figure}[H]
  \centering
  \includegraphics[width=0.94\textwidth]{figures/open_problem_crop.png}
  \caption{The beta-Ginibre question in arXiv:1805.00708, PDF page 8.}
\end{figure}

We solve the Poincar\'e part exactly when $n=2$.  It is useful to keep the
quadratic confinement coefficient general.  Let $\mu_{a,\beta}$ be the
probability measure on $\C^2\cong\R^4$ with density
\begin{equation}\label{eq:measure}
 \frac{d\mu_{a,\beta}}{dz_1\,dz_2}
 =\frac1{Z_{a,\beta}}
 e^{-a(|z_1|^2+|z_2|^2)}|z_1-z_2|^\beta,
 \qquad a>0,\quad\beta\geq0.
\end{equation}
Gradients below are real Euclidean gradients on the two copies of $\C$.

\begin{theorem}[Exact two-particle constant]\label{thm:main}
For every real $f$ in the weighted Sobolev space $H^1(\mu_{a,\beta})$,
\begin{equation}\label{eq:PI}
 \Var_{\mu_{a,\beta}}(f)
 \leq \frac{\sqrt{\beta^2+4}+\beta}{4a}
 \int_{\C^2}\bigl(|\nabla_{z_1}f|^2+|\nabla_{z_2}f|^2\bigr)\dmu.
\end{equation}
The constant is optimal.  Equivalently, the spectral gap is
\begin{equation}\label{eq:gap}
 \lambda_{a,\beta}=a\bigl(\sqrt{\beta^2+4}-\beta\bigr).
\end{equation}
If $\beta>0$, put
\[
 \alpha_1=\frac{\sqrt{\beta^2+4}-\beta}{2}.
\]
Apart from additive constants, equality is attained precisely by the real
span of
\[
 |z_1-z_2|^{\alpha_1-1}\Re(z_1-z_2),\qquad
 |z_1-z_2|^{\alpha_1-1}\Im(z_1-z_2).
\]
For $\beta=0$, all affine real functions attain equality.
\end{theorem}

\begin{corollary}[Source normalization]\label{cor:source}
For the two-particle beta-Ginibre law in arXiv:1805.00708, $a=n=2$, so
\[
 C_{\mathrm P}=\frac{\sqrt{\beta^2+4}+\beta}{8},\qquad
 \lambda=2\bigl(\sqrt{\beta^2+4}-\beta\bigr).
\]
For the classical complex Ginibre exponent $\beta=2$ this gives
\[
 C_{\mathrm P}=\frac{1+\sqrt2}{4},\qquad
 \lambda=4(\sqrt2-1).
\]
\end{corollary}

\section{Proof intuition}

Use orthonormal center and difference coordinates
\[
 u=\frac{z_1+z_2}{\sqrt2},\qquad
 v=\frac{z_1-z_2}{\sqrt2}.
\]
The density then factors.  The center $u$ is an ordinary planar Gaussian,
while the relative coordinate $v$ has density
$|v|^\beta e^{-a|v|^2}$.  The only non-Gaussian issue is therefore a
two-dimensional radial oscillator with a singular drift at the origin.

In polar coordinates, separate the angular mode $e^{im\theta}$.  A trial
factor $r^\alpha$ cancels the inverse-square term exactly when
$\alpha(\alpha+\beta)=m^2$.  After the change $\rho=ar^2$, the remaining
equation is the generalized Laguerre equation.  This gives the whole
spectrum explicitly.  The first excitation is not radial: it is the
$|m|=1$ angular mode, and its noninteger radial power is what makes the
optimal constant larger than the Gaussian center-of-mass constant.

\section{Factorization}

The map $(z_1,z_2)\mapsto(u,v)$ is orthogonal on $\R^4$, and
\[
 |z_1|^2+|z_2|^2=|u|^2+|v|^2,\qquad
 |z_1-z_2|^\beta=2^{\beta/2}|v|^\beta.
\]
The harmless power of $2$ is absorbed by normalization.  Hence
\begin{equation}\label{eq:product}
 \mu_{a,\beta}=\gamma_a\otimes\nu_{a,\beta},
\end{equation}
where
\[
 d\gamma_a(u)=\frac a\pi e^{-a|u|^2}\,du,
 \qquad
 d\nu_{a,\beta}(v)=c_{a,\beta}|v|^\beta e^{-a|v|^2}\,dv.
\]
The coordinate change also preserves the squared gradient.  The planar
Gaussian factor has spectral gap $2a$.  It remains to compute the gap of
$\nu_{a,\beta}$.

\section{The complete relative spectrum}

In polar coordinates $v=(r,\theta)$, the relative measure is proportional to
\[
 r^{\beta+1}e^{-ar^2}\,dr\,d\theta.
\]
The self-adjoint generator associated with the closure of the Dirichlet form
$\int|\nabla f|^2d\nu_{a,\beta}$ is the Friedrichs realization of
\begin{equation}\label{eq:generator}
 \cL=\partial_r^2+\left(\frac{\beta+1}{r}-2ar\right)\partial_r
       +\frac1{r^2}\partial_\theta^2.
\end{equation}
Initially this identity may be read on smooth functions supported away from
$r=0$; the form closure supplies the natural boundary condition there.

For $m\in\mathbb Z$, define
\begin{equation}\label{eq:alpha}
 \alpha_m=\frac{\sqrt{\beta^2+4m^2}-\beta}{2},\qquad
 q_m=\alpha_m+\frac\beta2.
\end{equation}
Thus $\alpha_m\geq0$ and $\alpha_m(\alpha_m+\beta)=m^2$.

\begin{proposition}[Relative spectral decomposition]\label{prop:spectrum}
For $m\in\mathbb Z$ and $k\in\mathbb N_0$, let
\begin{equation}\label{eq:eigenfunction}
 \Phi_{m,k}(r,\theta)
 =r^{\alpha_m}L_k^{(q_m)}(ar^2)e^{im\theta},
\end{equation}
where $L_k^{(q)}$ is the generalized Laguerre polynomial.  These functions
form a complete orthogonal eigenfamily in the complexification of
$L^2(\nu_{a,\beta})$, and
\begin{equation}\label{eq:eigenvalues}
 -\cL\Phi_{m,k}=(4ak+2a\alpha_m)\Phi_{m,k}.
\end{equation}
Consequently
\[
 \operatorname{gap}(\nu_{a,\beta})
 =a\bigl(\sqrt{\beta^2+4}-\beta\bigr).
\]
\end{proposition}

\begin{proof}
Fix an angular mode $e^{im\theta}$ and write
$R(r)=r^\alpha y(\rho)$ with $\rho=ar^2$.  The radial operator in that mode is
\[
 \cL_mR=R''+\left(\frac{\beta+1}{r}-2ar\right)R'
           -\frac{m^2}{r^2}R.
\]
A direct differentiation, using $\alpha(\alpha+\beta)=m^2$, gives
\begin{equation}\label{eq:calculation}
 \cL_m\bigl(r^\alpha y(ar^2)\bigr)
 =4a r^\alpha\left[
 \rho y''+(\alpha+\tfrac\beta2+1-\rho)y'-\frac\alpha2y
 \right].
\end{equation}
The Laguerre equation
\[
 \rho y''+(q+1-\rho)y'+ky=0
\]
with $q=\alpha+\beta/2$ now proves
\eqref{eq:eigenvalues}.

For completeness, Fourier series decompose the angular variable.  In the
$m$th sector, multiplication by $r^{\alpha_m}$ and the substitution
$\rho=ar^2$ identify the radial $L^2$ norm, up to a constant, with
\[
 \int_0^\infty |y(\rho)|^2\rho^{q_m}e^{-\rho}\,d\rho.
\]
Since $q_m>-1$, the generalized Laguerre polynomials are complete in this
space.  The positive root $\alpha_m$ is the finite-energy root selected by
the Friedrichs form.  The other root is inadmissible: for $m\ne0$ it is not
square integrable, and for $m=0$ its nonconstant branch has infinite energy.
Thus the displayed family exhausts the spectrum.

The constant is $m=0,k=0$.  The nonzero radial eigenvalues begin at $4a$.
For $m\ne0$, $\alpha_m$ increases with $|m|$, so the first angular
eigenvalue is
\[
 2a\alpha_1=a\bigl(\sqrt{\beta^2+4}-\beta\bigr).
\]
Because $\sqrt{\beta^2+4}-\beta\leq2$, this is at most $2a$, and hence is
strictly below $4a$.  It is the relative spectral gap.
\end{proof}

\section{Proof of Theorem \ref{thm:main}}

By \eqref{eq:product}, the full generator is the tensor sum of the center
generator and the relative generator.  Its spectrum consists of sums of
factor eigenvalues, so its gap is the smaller of the factor gaps:
\[
 \min\left\{2a,\,a(\sqrt{\beta^2+4}-\beta)\right\}
 =a(\sqrt{\beta^2+4}-\beta).
\]
The reciprocal of this gap is
\[
 \frac1{a(\sqrt{\beta^2+4}-\beta)}
 =\frac{\sqrt{\beta^2+4}+\beta}{4a},
\]
which proves \eqref{eq:PI} with its optimal constant.

For $\beta>0$, the gap is strictly smaller than the center gap.  The real
$m=\pm1,k=0$ eigenfunctions in the relative coordinate are
$r^{\alpha_1}\cos\theta$ and $r^{\alpha_1}\sin\theta$.  Rewriting them in
terms of $z_1-z_2$ gives the equality cases stated in the theorem.  They lie
in the form domain: near zero their energy density is integrable, and cutoff
approximations have vanishing extra energy.  When $\beta=0$, both factors are
ordinary planar Gaussians with the same gap $2a$, and all affine real
functions are extremal.  This completes the proof.

\section{Checks, scope, and novelty status}

The formula has the required Gaussian limit: at $\beta=0$ it gives gap
$2a$.  It also has the strong-repulsion asymptotic
\[
 \lambda_{a,\beta}\sim\frac{2a}{\beta},\qquad
 C_{\mathrm P}\sim\frac{\beta}{2a}.
\]
An accompanying exact symbolic script verifies the radial conjugation
identity, the Laguerre residual, the ordering of the candidate eigenvalues,
and the specializations $\beta=0$ and $\beta=2$.

This is a complete answer for the Poincar\'e constant of the two-particle
family, hence a solved subcase of the source's arbitrary-$n$ question.  It
does not determine the constant for arbitrary particle number or the
logarithmic Sobolev constant.  A bounded novelty search through the run
indexes, exact-title searches, beta-Ginibre functional-inequality searches,
and searches for the resulting closed formula found no prior statement of
this exact two-particle constant.  Chafa\"i, Bolley, and Fontbona prove
existence of a Poincar\'e inequality and discuss non-explicit bounds for the
planar Coulomb gas \cite{CBF}; that result is consistent with, but does not
contain, the spectral calculation above.  This novelty check is not an
exhaustive bibliographic certification.

For human review, the two points deserving the closest check are the
Friedrichs-domain selection at the collision point and the claim that the
Fourier--Laguerre family exhausts the relative spectrum.  Both are addressed
in Proposition~\ref{prop:spectrum}; an independent symbolic residual and
Rayleigh-quotient check is included with the packet.

\begin{thebibliography}{9}
\bibitem{ChafaiLehec}
D.~Chafa\"i and J.~Lehec,
\emph{On Poincar\'e and logarithmic Sobolev inequalities for a class of
singular Gibbs measures}, arXiv:1805.00708 (2018),
\url{https://arxiv.org/abs/1805.00708}.

\bibitem{CBF}
F.~Bolley, D.~Chafa\"i, and J.~Fontbona,
\emph{Dynamics of a planar Coulomb gas}, arXiv:1706.08776 (2017),
\url{https://arxiv.org/abs/1706.08776}.
\end{thebibliography}

\end{document}
